Integrating large language models (LLMs) with existing software requires more than selecting suitable tools: multi-step actions must be coordinated, dialogue context must be maintained, and consequential operations must remain under human control.
This paper presents the design, configuration selection and evaluation of an MCP-based conversational scheduling assistant.
The prototype exposes existing service capabilities through the Model Context Protocol (MCP) and combines structured dialogue state, server-side validation and explicit user confirmation before booking.
A selection study compares three models and three orchestration variants on two scenarios, yielding 90 analysed runs after excluding cold starts.
Claude Haiku 4.5 with individual tools and no separate planner produced correct outcomes in all five repetitions of each scenario and was selected for further evaluation.
The comparison illustrates that technically successful tool calls need not produce correct outcomes.
After a separate settings experiment with 50 analysed runs, only the chosen Haiku configuration was assessed on seven functional cases, 80 executions of eight direct prompt-injection inputs, concurrent use and exploratory feedback from three operational users and six developers.
The functional cases met expectations after an error-handling adjustment for empty search results.
No control violation was observed in the injection tests, although some inputs disrupted the expected workflow.
User feedback favoured flexible input while identifying a need for concise responses and shorter delays.
Selection results and subsequent observations provide task-specific evidence, not an independent model ranking or proof of production readiness.
